\chapter{Kinetic Theory and the Perfect Gas}\label{ch:b1:kinetic-theory}

The air in a classroom weighs as much as a grown man, and contains more
molecules than there are grains of sand on Earth, each flying at the
speed of a rifle bullet and colliding ten billion times a second with
its neighbors. Nobody can follow one of them; nobody needs to. A handful
of averages --- \omterm{def:b1:kinetic-theory:pressure}{pressure}, \omterm{def:b1:kinetic-theory:temperature}{temperature}, volume, amount --- describe the
air completely for every purpose of engineering, and a simple picture
of molecules bouncing off walls explains where those averages come
from and how they are related. This chapter opens the thermodynamics
part of the volume by looking at matter from both ends: the microscopic
chaos and the macroscopic calm, and the kinetic theory that connects
them in the case of the \omterm{def:b1:kinetic-theory:model}{perfect gas}.

\section{Two scales of description}

\begin{definition}[Microscopic and macroscopic scales]\label{def:b1:kinetic-theory:scales}
\index{microscopic scale}\index{macroscopic scale}\index{Avogadro constant}\index{mole}
At the \emph{microscopic} scale matter is made of molecules: $\qty{1}{mol}$
contains $N_A = \qty{6.02e23}{}$ of them, each about \qty{0.3}{nm} across;
in a gas at ordinary conditions they are \qty{3}{nm} apart, move at
hundreds of meters per second and collide every \qty{70}{nm} or so.
At the \emph{macroscopic} scale a system is described by a few
\emph{\omterm{def:b1:kinetic-theory:equilibrium}{state variables}} --- \omterm{def:b1:kinetic-theory:pressure}{pressure} $P$, volume $V$, \omterm{def:b1:kinetic-theory:temperature}{temperature} $T$,
amount of substance $n$ --- that are averages over enormous numbers of
molecules. A variable is \emph{extensive} if it doubles when the system
is doubled ($V$, $n$, mass, energy), \emph{intensive} if it does not
($P$, $T$, density).
\end{definition}

\begin{definition}[Thermodynamic equilibrium]\label{def:b1:kinetic-theory:equilibrium}
\index{thermodynamic equilibrium}\index{state variable}
A system is in \emph{thermodynamic equilibrium} when its state variables
are uniform and constant in time and no macroscopic flow of matter or
energy crosses it. Only then are $P$ and $T$ defined for the system as a
whole; the \emph{equation of state} links them to $V$ and $n$.
\end{definition}

\begin{definition}[Pressure]\label{def:b1:kinetic-theory:pressure}
\index{pressure}\index{pascal}\index{bar}
A fluid at rest pushes on every surface element $\dd S$ of a wall with a
force $\dd\vect F = P\,\dd S\,\vect n$ normal to it, toward the wall; $P$ is
the \emph{pressure}, in pascals ($\qty{1}{Pa} = \qty{1}{N/m^2}$); $\qty{1}{bar}
= \qty{1e5}{Pa}$, $\qty{1}{atm} = \qty{1.013e5}{Pa}$. At a given point it does
not depend on the orientation of the surface (\cref{ch:b1:fluid-statics}).
\end{definition}

\begin{definition}[Temperature]\label{def:b1:kinetic-theory:temperature}
\index{temperature}\index{kelvin}\index{zeroth law}
Two systems in contact through a wall that lets energy pass reach, in
time, a common state: they are then at the same \emph{temperature}, and
two bodies each in equilibrium with a third are in equilibrium with
each other (the \emph{zeroth law}). The kelvin scale is fixed by
$k_B = \qty{1.380649e-23}{J/K}$ exactly; Celsius: $\theta = T - 273.15$. The
kinetic theory below gives $T$ its microscopic meaning.
\end{definition}

\section{The kinetic model of the perfect gas}

\begin{definition}[Perfect gas, microscopic model]\label{def:b1:kinetic-theory:model}
\index{perfect gas}\index{ideal gas}
A \emph{perfect gas} (ideal gas) is a collection of $N$ molecules,
treated as points of mass $m$, moving freely and at random --- no
interaction except brief collisions --- in a container of volume $V$.
At equilibrium the velocity distribution is isotropic and stationary;
$n^* = N/V$ is the number density and $\langle v^2\rangle$ the mean square
speed.
\end{definition}

\begin{theorem}[Kinetic pressure]\label{thm:b1:kinetic-theory:pressure}
\index{kinetic pressure}
The \omterm{def:b1:kinetic-theory:pressure}{pressure} exerted by the gas on the walls is
\[
P = \tfrac13\,n^*m\langle v^2\rangle = \tfrac23\,n^*\langle\epsilon_k\rangle ,
\]
where $\langle\epsilon_k\rangle = \tfrac12 m\langle v^2\rangle$ is the mean
\omterm{thm:b1:work-and-energy:ke}{kinetic energy} of a molecule.
\end{theorem}

\begin{proof}[Partial proof]
Take the simplified model in which the molecules all have speed $u$ and
move along the three axes, one sixth in each direction. A molecule
hitting the wall $x = 0$ head-on bounces back elastically and delivers
the momentum $2mu$ to it. In time $\dd t$, the molecules moving toward
the wall that reach an area $S$ are those within the slab $u\,\dd t$
thick, one sixth of the $n^*Su\,\dd t$ molecules there. Force $= $
momentum per time $= (n^*Su\,\dd t/6)(2mu)/\dd t = \tfrac13 n^*mu^2S$, so
$P = \tfrac13 n^*mu^2$. Averaging over the actual distribution of speeds
and directions replaces $u^2$ by $\langle v^2\rangle$ (the isotropy gives
$\langle v_x^2\rangle = \tfrac13\langle v^2\rangle$); the Year 2 volume does
it in full.
\end{proof}

\begin{omfigure}
\begin{tikzpicture}[font=\small, x=1cm, y=1cm]
  \draw[thick] (0,-1.6) -- (0,1.6); \fill[pattern=north east lines] (-0.2,-1.6) rectangle (0,1.6);
  \draw[dashed] (1.8,-1.6) -- (1.8,1.6);
  \draw[Stealth-Stealth] (0,-1.85) -- (1.8,-1.85) node[midway, below] {$u\,\dd t$};
  \foreach \p in {(0.6,1.1),(1.2,0.5),(0.9,-0.2),(1.5,-0.9),(0.4,-1.2),(1.3,1.3)} {\fill[omDef] \p circle (2pt);}
  \draw[-Stealth, omDef, thick] (1.2,0.5) -- (0.5,0.5);
  \draw[-Stealth, omThm, thick] (0.25,0.2) -- (0.95,0.2);
  \node[omDef, font=\footnotesize, above] at (0.85,0.55) {$m\vect u$};
  \node[omThm, font=\footnotesize, below] at (0.6,0.15) {$-m\vect u$};
  \node[right, font=\footnotesize] at (2.0,0.9) {momentum $2mu$ given};
  \node[right, font=\footnotesize] at (2.0,0.55) {to the wall per impact};
  \node[right, font=\footnotesize] at (2.0,-0.5) {$\tfrac16 n^*Su\,\dd t$ impacts in $\dd t$};
  \begin{scope}[xshift=7.2cm, yshift=-1.6cm]
  \begin{axis}[omaxis, axis x line=bottom, axis y line=left, width=6.4cm, height=4.2cm, xlabel={$v$ (\unit{m/s})}, ylabel={$f(v)$},
      xlabel style={at={(axis description cs:0.5,-0.1)}, anchor=north},
      xmin=0, xmax=1400, ymin=0, ymax=1.2, xtick={500,1000}, ytick=\empty, anchor=origin, at={(0,0)}]
    \addplot[omDef, thick, domain=0:1400, samples=120] {(x/422)^2*exp(1-(x/422)^2)};
    \addplot[omThm, thick, domain=0:1400, samples=120] {0.707*(x/597)^2*exp(1-(x/597)^2)};
    \node[omDef, font=\footnotesize] at (axis cs:430,1.12) {\qty{300}{K}};
    \node[omThm, font=\footnotesize] at (axis cs:900,0.72) {\qty{600}{K}};
  \end{axis}
  \end{scope}
\end{tikzpicture}

{\small Left: the kinetic origin of \omterm{def:b1:kinetic-theory:pressure}{pressure} --- molecules striking a
wall and bouncing back deliver momentum to it; their number per second
times $2mu$ is the force. Right: the actual speeds are spread (Maxwell's
distribution, here for nitrogen at two \omterm{def:b1:kinetic-theory:temperature}{temperatures}); the model's single
$u$ is replaced by the \omterm{def:b1:kinetic-theory:kinetictemp}{root-mean-square speed}.}
\end{omfigure}

\begin{definition}[Kinetic temperature; equation of state]\label{def:b1:kinetic-theory:kinetictemp}
\index{kinetic temperature}\index{equation of state!perfect gas}\index{gas constant}\index{root-mean-square speed}
The \omterm{def:b1:kinetic-theory:temperature}{temperature} of a \omterm{def:b1:kinetic-theory:model}{perfect gas} is \emph{defined} by the mean \omterm{thm:b1:work-and-energy:ke}{kinetic
energy} of its molecules,
\[
\langle\epsilon_k\rangle = \tfrac12 m\langle v^2\rangle = \tfrac32 k_BT ,
\]
so that, combining with \cref{thm:b1:kinetic-theory:pressure},
\[
PV = Nk_BT = nRT, \qquad R = N_Ak_B = \qty{8.314}{J/(mol.K)} ,
\]
the \emph{equation of state of the \omterm{def:b1:kinetic-theory:model}{perfect gas}}. The \emph{root-mean-square
speed} is $u = \sqrt{\langle v^2\rangle} = \sqrt{3k_BT/m} = \sqrt{3RT/M}$,
$M$ the molar mass.
\end{definition}

\begin{example}[Numbers for air]\label{ex:b1:kinetic-theory:air}
Nitrogen at \qty{300}{K}: $u = \sqrt{3 \times 8.314 \times 300/0.028} =
\qty{517}{m/s}$; hydrogen, $\qty{1930}{m/s}$; the heavier the slower, as
$1/\sqrt M$. At \qty{1}{\omterm{def:b1:kinetic-theory:pressure}{bar}} and \qty{300}{K}, $n^* = P/k_BT = \qty{2.4e25}{m^{-3}}$:
\qty{24}{L} per \omterm{def:b1:kinetic-theory:scales}{mole}, \qty{2.7e22}{} molecules in a liter. A classroom of
\qty{200}{m^3}: \qty{8300}{mol}, \qty{240}{kg} of air.
\end{example}

\begin{remark}[Why the model works, and when it fails]\label{rem:b1:kinetic-theory:validity}
\index{mean free path}
Between collisions a molecule of air flies a \emph{\omterm{rem:b1:kinetic-theory:validity}{mean free path}}
$\lambda \approx 1/(\sqrt2\pi d^2n^*) \approx \qty{70}{nm}$ --- two hundred
times its size: the gas is mostly empty, interactions are rare, and
$PV = nRT$ holds to better than a percent at ordinary \omterm{def:b1:kinetic-theory:pressure}{pressures}. At high
\omterm{def:b1:kinetic-theory:pressure}{pressure} or low \omterm{def:b1:kinetic-theory:temperature}{temperature} the molecules' own volume and their
attraction matter: van der Waals's equation $(P + an^2/V^2)(V - nb) = nRT$
corrects for both, and below a critical \omterm{def:b1:kinetic-theory:temperature}{temperature} the gas liquefies
(\cref{ch:b1:phase-changes}).
\end{remark}

\begin{omfigure}
\begin{tikzpicture}
\begin{axis}[omaxis, axis x line=bottom, axis y line=left, width=10cm, height=6cm, xlabel={$V$}, ylabel={$P$},
    xlabel style={at={(axis description cs:0.5,-0.08)}, anchor=north},
    xmin=0, xmax=5.2, ymin=-0.4, ymax=4.2, xtick=\empty, ytick=\empty]
  \addplot[omDef, thick, domain=0.3:5.2, samples=100] {1.0/x};
  \addplot[omDef, thick, domain=0.45:5.2, samples=100] {1.5/x};
  \addplot[omDef, thick, domain=0.6:5.2, samples=100] {2.0/x};
  \addplot[omThm, thick, domain=0.75:5.2, samples=100] {2.5/x};
  \node[omDef, font=\footnotesize, below] at (axis cs:4.9,0.19) {$T_1$};
  \node[omThm, font=\footnotesize, above] at (axis cs:4.9,0.53) {$T_4$};
  \draw[-Stealth, omExo] (axis cs:1.3,1.0) -- (axis cs:2.3,2.0) node[above right, font=\footnotesize] {$T$ increasing};
  \node[font=\footnotesize] at (axis cs:3.4,3.2) {$PV = nRT$: isotherms $T_1 < \dots < T_4$};
\end{axis}
\end{tikzpicture}

{\small The Clapeyron diagram of a \omterm{def:b1:kinetic-theory:model}{perfect gas}: at each \omterm{def:b1:kinetic-theory:temperature}{temperature} the
isotherm is a hyperbola $P = nRT/V$, higher for higher $T$. A transformation
is a path in this plane; a cycle, a closed loop (\cref{ch:b1:heat-engines}).}
\end{omfigure}

\section{Internal energy of the perfect gas}

\begin{definition}[Internal energy]\label{def:b1:kinetic-theory:U}
\index{internal energy}
The \emph{internal energy} $U$ of a system is the sum of the kinetic
energies of its molecules in the frame where the system is at rest and
of their mutual potential energies; it is an extensive state function.
\end{definition}

\begin{proposition}[Internal energy of perfect gases]\label{prop:b1:kinetic-theory:Ugas}
\index{monatomic gas}\index{diatomic gas}\index{heat capacity at constant volume}\index{Joule's laws}
For a \omterm{def:b1:kinetic-theory:model}{perfect gas} $U$ depends on $T$ alone (first Joule law):
\begin{gather*}
U = \tfrac32 nRT \quad (\text{monatomic: He, Ne, Ar}), \\
U = \tfrac52 nRT \quad (\text{diatomic at ordinary temperatures: N}_2, \text{O}_2, \text{H}_2),
\end{gather*}
so the molar \omterm{prop:b1:kinetic-theory:Ugas}{heat capacity at constant volume} $C_{V,m} = \dd U_m/\dd T$
is $\tfrac32 R = \qty{12.5}{J/(mol.K)}$ and $\tfrac52 R = \qty{20.8}{J/(mol.K)}$
respectively.
\end{proposition}

\begin{proof}[Partial proof]
No mutual \omterm{def:b1:work-and-energy:conservative}{potential energy} by hypothesis, so $U = N\langle\epsilon_k\rangle
= \tfrac32 Nk_BT$ for point molecules. A diatomic molecule also rotates:
each of its two rotational degrees of freedom carries, at equilibrium,
the same $\tfrac12 k_BT$ as each translational one (the \emph{equipartition}
theorem, admitted; the Year 3 volume derives it and explains why the
vibration joins in only above a thousand \omterm{def:b1:kinetic-theory:temperature}{kelvins}).
\end{proof}

\begin{example}[A room's worth]\label{ex:b1:kinetic-theory:room}
The \qty{8300}{mol} of air in the classroom at \qty{300}{K}: $U = \tfrac52 \times
8300 \times 8.314 \times 300 = \qty{5.2e7}{J}$ --- the energy of a liter and a
half of gasoline, held in molecular motion. Warming the room by $1$ K
costs $\tfrac52 nR = \qty{170}{kJ}$ (at constant volume).
\end{example}

\section{Condensed phases}

\begin{proposition}[Model of the incompressible, indilatable phase]\label{prop:b1:kinetic-theory:condensed}
\index{condensed phase}\index{thermal expansion}\index{compressibility}
Liquids and solids are modeled, at this level, as phases of fixed
volume --- neither compressible nor dilatable --- whose \omterm{def:b1:kinetic-theory:U}{internal energy}
depends on $T$ alone: $\dd U = C\,\dd T$ with $C = mc$, $c$ the specific
heat capacity (\qty{4180}{J/(kg.K)} for water, \qty{450}{J/(kg.K)} for
steel). The real, small deviations are measured by the \omterm{prop:b1:kinetic-theory:condensed}{thermal expansion}
coefficient $\alpha = (1/V)(\partial V/\partial T)_P$ ($\sim\qty{1e-5}{K^{-1}}$
for metals, \qty{2e-4}{K^{-1}} for liquids) and the isothermal
compressibility $\chi_T = -(1/V)(\partial V/\partial P)_T$ ($\sim\qty{5e-10}{Pa^{-1}}$
for water, against $1/P \approx \qty{1e-5}{Pa^{-1}}$ for a gas).
\end{proposition}
\admitted

\begin{example}[How incompressible]\label{ex:b1:kinetic-theory:water}
Water at the bottom of the Mariana trench (\qty{1100}{\omterm{def:b1:kinetic-theory:pressure}{bar}}) is compressed
by $\chi_T\Delta P \approx 5\times10^{-10} \times 1.1\times10^8 = 5\%$; a
\qty{30}{m} steel rail warmed by $50$ K lengthens by $\alpha L\Delta T =
\qty{18}{mm}$ --- the reason for expansion joints. Against a gas, whose
volume halves under a doubled \omterm{def:b1:kinetic-theory:pressure}{pressure}, the model of a fixed volume is
good to a few percent in most situations of this volume.
\end{example}

\section{Exercises}

\begin{exercise}[$\star$]\label{exo:b1:kinetic-theory:1}
Number of molecules in \qty{1.0}{L} of gas at \qty{0}{\degreeCelsius} and
\qty{1}{atm}; mean distance between neighbors; compare with the molecular
size \qty{0.3}{nm}.
\end{exercise}

\begin{exercise}[$\star$]\label{exo:b1:kinetic-theory:2}
\omterm{def:b1:kinetic-theory:kinetictemp}{Root-mean-square speeds} of H$_2$, N$_2$, O$_2$ and CO$_2$ at \qty{300}{K},
and of N$_2$ at \qty{77}{K} (boiling nitrogen).
\end{exercise}

\begin{exercise}[$\star$]\label{exo:b1:kinetic-theory:3}
A car tire of volume \qty{30}{L} is inflated to a gauge \omterm{def:b1:kinetic-theory:pressure}{pressure} of
\qty{2.5}{\omterm{def:b1:kinetic-theory:pressure}{bar}} at \qty{20}{\degreeCelsius}. Amount and mass of air inside;
absolute \omterm{def:b1:kinetic-theory:pressure}{pressure} after a drive that heats it to \qty{50}{\degreeCelsius}.
\end{exercise}

\begin{exercise}[$\star$]\label{exo:b1:kinetic-theory:4}
Estimate the \omterm{rem:b1:kinetic-theory:validity}{mean free path} of air molecules at \qty{1}{\omterm{def:b1:kinetic-theory:pressure}{bar}}, \qty{300}{K}
($d = \qty{0.37}{nm}$), and the collision frequency of one molecule. At
what \omterm{def:b1:kinetic-theory:pressure}{pressure} does the \omterm{rem:b1:kinetic-theory:validity}{mean free path} reach \qty{1}{m} (a ``vacuum'')?
\end{exercise}

\begin{exercise}[$\star\star$]\label{exo:b1:kinetic-theory:5}
Density of air at \qty{1.013}{\omterm{def:b1:kinetic-theory:pressure}{bar}} and \qty{20}{\degreeCelsius}
($M = \qty{29}{g/mol}$), then at \qty{100}{\degreeCelsius}. Lift of a
\qty{2800}{m^3} hot-air balloon (difference of the two weights of air).
\end{exercise}

\begin{exercise}[$\star\star$]\label{exo:b1:kinetic-theory:6}
Air is $21\%$ O$_2$ by molecules. Partial \omterm{def:b1:kinetic-theory:pressure}{pressure} of oxygen at sea level;
at \qty{5000}{m} where $P = \qty{0.54}{atm}$; at the summit of Everest
(\qty{0.33}{atm}). Why does the body struggle there?
\end{exercise}

\begin{exercise}[$\star\star$]\label{exo:b1:kinetic-theory:7}
One \omterm{def:b1:kinetic-theory:scales}{mole} of CO$_2$ at \qty{300}{K} in \qty{1.0}{L}: \omterm{def:b1:kinetic-theory:pressure}{pressure} from the
perfect-gas law, then from van der Waals ($a = \qty{0.364}{Pa.m^6/mol^2}$,
$b = \qty{4.27e-5}{m^3/mol}$). Which correction dominates here?
\end{exercise}

\begin{exercise}[$\star\star$]\label{exo:b1:kinetic-theory:8}
\omterm{def:b1:kinetic-theory:U}{Internal energy} of the air in a \qty{50}{m^3} room at \qty{1}{\omterm{def:b1:kinetic-theory:pressure}{bar}},
\qty{300}{K}; energy to warm it by \qty{5}{K} at constant volume; time
with a \qty{2}{kW} heater (no losses).
\end{exercise}

\begin{exercise}[$\star\star$]\label{exo:b1:kinetic-theory:9}
A \qty{30}{m} steel rail ($\alpha = \qty{1.2e-5}{K^{-1}}$) between
\qty{-10}{\degreeCelsius} and \qty{40}{\degreeCelsius}: length change.
A brass ring ($\alpha = \qty{1.9e-5}{K^{-1}}$) of inner diameter
\qty{49.95}{mm} must slip over a \qty{50.00}{mm} shaft: by how much must
it be heated?
\end{exercise}

\begin{exercise}[$\star\star\star$]\label{exo:b1:kinetic-theory:10}
Energy per molecule at \qty{300}{K} ($\tfrac32 k_BT$) in joules and eV;
molar heat capacities $C_{V,m}$ of argon and of nitrogen; why is that of
nitrogen larger though its molecules are lighter?
\end{exercise}

\begin{exercise}[$\star\star\star$]\label{exo:b1:kinetic-theory:11}
Redo the derivation of the \omterm{thm:b1:kinetic-theory:pressure}{kinetic pressure} in the six-direction model,
then show that for an isotropic distribution of velocities $\langle v_x^2\rangle
= \tfrac13\langle v^2\rangle$ and that the \omterm{def:b1:kinetic-theory:pressure}{pressure} formula is unchanged.
\end{exercise}

\begin{exercise}[$\star\star\star$]\label{exo:b1:kinetic-theory:12}
\omterm{def:b1:kinetic-theory:kinetictemp}{Root-mean-square speeds} of H$_2$ and O$_2$ at \qty{300}{K} compared with
the Earth's \omterm{prop:b1:central-forces:escape}{escape velocity}, and with the Moon's (\qty{2.4}{km/s}).
Knowing that the fraction of molecules faster than $3u$ is small but not
zero, explain why the Earth has lost its hydrogen and the Moon its whole
atmosphere.
\end{exercise}

\section{Problem: The air in the room}

\begin{problem}[{Weekend problem --- a classroom of air: counting its
molecules, weighing it, timing their flights, and recovering from the
same model the lift of a balloon, the thinness of mountain air and the
speed of sound}]
\label{pb:b1:kinetic-theory:1}
A classroom measures $10 \times 8 \times \qty{2.5}{m}$; its air is at
$P = \qty{1.013e5}{Pa}$, $T = \qty{293}{K}$, molar mass $M = \qty{29.0}{g/mol}$,
$79\%$ N$_2$ and $21\%$ O$_2$ by molecules. $R = \qty{8.314}{J/(mol.K)}$,
$k_B = \qty{1.38e-23}{J/K}$, $N_A = \num{6.02e23}$, $g = \qty{9.81}{m/s^2}$.

\medskip
\noindent\textbf{Part I --- Counting and weighing.}
\begin{enumerate}
  \item Amount of air (\omterm{def:b1:kinetic-theory:scales}{moles}) and number of molecules in the room.
  \item Mass of the air; compare with a person.
  \item Number density $n^*$ and mean distance between molecules
        ($n^{*-1/3}$); compare with the molecular diameter \qty{0.37}{nm}.
  \item \omterm{def:b1:kinetic-theory:kinetictemp}{Root-mean-square speed} of N$_2$ and of O$_2$; why does the heavier
        gas move more slowly at the same \omterm{def:b1:kinetic-theory:temperature}{temperature}?
  \item Mean \omterm{thm:b1:work-and-energy:ke}{kinetic energy} of one molecule; total translational \omterm{thm:b1:work-and-energy:ke}{kinetic
        energy} of the room's air; compare with the \omterm{thm:b1:work-and-energy:ke}{kinetic energy} of a
        car at \qty{100}{km/h}.
  \item \omterm{rem:b1:kinetic-theory:validity}{Mean free path} $\lambda \approx 1/(\sqrt2\pi d^2n^*)$ and the mean
        time between collisions of a molecule.
\end{enumerate}

\medskip
\noindent\textbf{Part II --- \omterm{def:b1:kinetic-theory:pressure}{Pressure} from impacts.}
\begin{enumerate}[resume]
  \item Using the six-direction model with $u = \qty{500}{m/s}$, compute
        the number of impacts per second on \qty{1}{cm^2} of wall.
  \item Compute the momentum delivered per second, and check that it
        reproduces the atmospheric \omterm{def:b1:kinetic-theory:pressure}{pressure}.
  \item Total force of the air on one wall of \qty{8}{m} by \qty{2.5}{m};
        why does the wall not move?
  \item The room is warmed to \qty{303}{K} with its windows closed and
        sealed: new \omterm{def:b1:kinetic-theory:pressure}{pressure}, and the net force on the wall now.
  \item Same warming with a window open: what leaves, and how much?
\end{enumerate}

\medskip
\noindent\textbf{Part III --- Buoyancy.}
\begin{enumerate}[resume]
  \item Density of the room's air, and of air at \qty{373}{K} at the same
        \omterm{def:b1:kinetic-theory:pressure}{pressure}.
  \item A hot-air balloon of volume \qty{2800}{m^3} is filled with air at
        \qty{373}{K}: weight of the hot air, weight of the displaced cold
        air, and the available lift (Archimedes,
        \cref{ch:b1:fluid-statics}).
  \item The envelope, basket, burner and gas weigh \qty{300}{kg}: how many
        \qty{75}{kg} passengers can it lift?
  \item To what \omterm{def:b1:kinetic-theory:temperature}{temperature} should the air be heated to double the lift,
        and why is that not done (the envelope is nylon)?
  \item A helium balloon of the same volume at \qty{293}{K}
        ($M = \qty{4}{g/mol}$): lift? Why do tourist balloons use hot air
        nonetheless?
\end{enumerate}

\medskip
\noindent\textbf{Part IV --- Thin air and the speed of sound.}
\begin{enumerate}[resume]
  \item Partial \omterm{def:b1:kinetic-theory:pressure}{pressure} of oxygen in the room.
  \item At the summit of Everest $P = \qty{0.33}{atm}$ and $T \approx \qty{250}{K}$:
        number density of oxygen molecules compared with the room's.
  \item The atmosphere's \omterm{def:b1:kinetic-theory:pressure}{pressure} falls with height roughly as
        $P(z) = P_0\eu^{-z/H}$ with $H = RT/Mg$ (\cref{ch:b1:fluid-statics}):
        compute $H$ for \qty{260}{K} and the height at which $P = P_0/2$.
  \item The speed of sound in a \omterm{def:b1:kinetic-theory:model}{perfect gas} is $c_s = \sqrt{\gamma RT/M}$
        with $\gamma = 1.4$ for air (\cref{ch:b1:first-law}): compute it at
        \qty{293}{K} and compare with the rms speed; interpret the
        closeness.
  \item How does $c_s$ vary with $T$? Speed of sound in the room at
        \qty{303}{K}, and in helium at \qty{293}{K} ($\gamma = 5/3$): why
        does a helium voice sound high?
  \item At the top of Everest, is sound faster or slower than in the
        room? By how much?
  \item The room's air is replaced by argon at the same $P$ and $T$:
        what changes in the count, the mass, the speeds, the \omterm{def:b1:kinetic-theory:U}{internal
        energy}?
  \item \omterm{def:b1:kinetic-theory:U}{Internal energy} of the room's air ($U = \tfrac52 nRT$) and the
        energy to warm it by \qty{10}{K} at constant volume.
  \item Summarize in three lines what the kinetic model delivered: which
        macroscopic quantities it explained, and from which single
        microscopic average.
\end{enumerate}
\end{problem}
